\documentclass[12pt,a4paper]{article}
\usepackage[T1]{fontenc}
\usepackage[utf8]{inputenc}
\usepackage{times}
\usepackage[a4paper,top=2.5cm,bottom=2.5cm,left=2.5cm,right=2.5cm]{geometry}
\usepackage{graphicx}\usepackage{booktabs}\usepackage{array}
\usepackage{caption}\usepackage{fancyhdr}\usepackage{parskip}
\usepackage[hidelinks]{hyperref}\usepackage{float}
\renewcommand{\topfraction}{0.9}\renewcommand{\bottomfraction}{0.7}
\renewcommand{\textfraction}{0.08}\renewcommand{\floatpagefraction}{0.75}
\captionsetup{font=small,labelfont=bf,labelsep=colon,
              justification=raggedright,singlelinecheck=false}
\pagestyle{fancy}\fancyhf{}
\fancyhead[L]{\footnotesize KOUAME KOFFI FID\`ELE}
\fancyhead[C]{\footnotesize EXCEL SALES REPORT}
\fancyhead[R]{\footnotesize Data Analysis Practice}
\fancyfoot[C]{\thepage}\renewcommand{\headrulewidth}{0.4pt}
\newcommand{\F}[1]{\texttt{\small #1}}
\begin{document}

\begin{titlepage}\centering
\vspace*{3cm}
{\large\scshape Data Analysis Practice}\\[2.6cm]
\rule{\linewidth}{0.5pt}\\[0.8cm]
{\LARGE\bfseries Excel Sales Report}\\[0.6cm]
{\large Seven Practical Tasks on a Car-Dealership Sales Log}\\[0.8cm]
\rule{\linewidth}{0.5pt}\\[1.2cm]
{\normalsize Formatting $\bullet$ Conditional Formatting $\bullet$ Lookups $\bullet$
Formulas $\bullet$ Invoicing}\\[3.2cm]
{\large\itshape Prepared by}\\[0.5cm]
{\Large\bfseries KOUAME KOFFI FID\`ELE}\\[0.6cm]
{\large September 14, 2026}
\vfill\end{titlepage}

\tableofcontents
\newpage

\section{Overview}

This report documents a single Excel workbook,
\F{KOUAME\_Koffi\_Fidele\_Excel\_Tasks.xlsx}, built around a 30-order car-sales log
for May 2019: order date, sales agent, car type, unit price and count. Seven
independent sheets, one per task, apply a different Excel skill to that same
underlying data: professional formatting, conditional formatting, target-based
highlighting with MAX/MIN, summary formulas, a per-salesperson report, a price
lookup, and a printable invoice.

Every result in the workbook is a live formula, never a typed-in number. Changing a
price, a count, or the target in cell J3 recalculates every total, every average and
every highlight automatically. The workbook was validated with a full recalculation
pass: 244 formulas, zero errors.

\textbf{Conventions used throughout.} A single dark-navy header bar identifies each
table; column headers state the currency once (EGP) rather than repeating it in every
cell; dates use a plain \F{dd/mm/yyyy} format; and every sheet is set to print on a
single landscape page.

\section{Task 1: Formatting}

The raw sales log is turned into a readable table: a merged title bar, a bold header
row, thousands-separated currency figures, consistent date formatting, borders around
every cell, light row banding for readability, and a totals row driven by \F{SUM}.

\begin{figure}[tbp]\centering
\includegraphics[width=0.78\linewidth]{fig1.png}
\caption{Task 1: the Sales Report after formatting: dates, currency, borders and a
totals row.}\end{figure}

\section{Task 2: Conditional Formatting}

The highest and lowest Order Amount (the Total Sales column) are identified and
highlighted automatically, using two formula-based conditional formatting rules
rather than manual colouring:

\begin{table}[htbp]\centering
\begin{tabular}{@{}ll@{}}\toprule
\textbf{Rule} & \textbf{Formula} \\ \midrule
Highest value, green & \F{=F6=MAX(\$F\$6:\$F\$35)} \\
Lowest value, red & \F{=F6=MIN(\$F\$6:\$F\$35)} \\
\bottomrule\end{tabular}\end{table}

Because the rule is formula-based, the highlight follows the data: if any order
amount changes, the green and red cells move with it.

\begin{figure}[tbp]\centering
\includegraphics[width=0.76\linewidth]{fig2.png}
\caption{Task 2: highest Order Amount in green, lowest in red, with a legend.}
\end{figure}

\section{Task 3: Conditional Formatting Against a Target, with MAX/MIN}

A monthly results table (nine sales people, January to June) is compared against a
single Target cell, \F{J3 = 1000}. Every result above the target is shaded green,
every result below it is shaded red, using a \F{CellIsRule} anchored on the absolute
reference \F{\$J\$3} so the same two rules cover the whole table. The highest and
lowest results across all nine people and six months are then pulled out with:

\begin{table}[htbp]\centering
\begin{tabular}{@{}ll@{}}\toprule
\textbf{Cell} & \textbf{Formula} \\ \midrule
Lowest Result (J4) & \F{=MIN(B6:G14)} \\
Highest Result (J5) & \F{=MAX(B6:G14)} \\
\bottomrule\end{tabular}\end{table}

\begin{figure}[tbp]\centering
\includegraphics[width=0.94\linewidth]{fig3.png}\\[6pt]
\includegraphics[width=0.94\linewidth]{fig4.png}
\caption{Task 3: results above the target in green, below it in red, with the Target,
Lowest and Highest Result called out on the right.}\end{figure}

\section{Task 4: Summary Formulas}

Four summary figures are calculated directly from the order table, with no manual
arithmetic:

\begin{table}[htbp]\centering
\begin{tabular}{@{}lll@{}}\toprule
\textbf{Item} & \textbf{Formula} & \textbf{Result} \\ \midrule
Total Sales Amount & \F{=SUM(F6:F35)} & EGP 29,365,000 \\
Total Sold Cars & \F{=SUM(E6:E35)} & 50 \\
Count of Orders & \F{=COUNTA(A6:A35)} & 30 \\
Average Order Amount & \F{=AVERAGE(F6:F35)} & EGP 978,833 \\
\bottomrule\end{tabular}\end{table}

\begin{figure}[tbp]\centering
\includegraphics[width=0.72\linewidth]{fig5.png}
\caption{Task 4: the four summary cells, filled with formulas rather than typed-in
figures.}\end{figure}

\section{Task 5: Sales Analysis}

A results report is built for each of the six sales agents in the log, using
\F{COUNTIF}, \F{SUMIF} and \F{AVERAGEIF} against the order table, and the Best Sales
Person is then identified with \F{INDEX}/\F{MATCH} on the highest Total Sales Amount:

\begin{table}[htbp]\centering\small
\begin{tabular}{@{}ll@{}}\toprule
\textbf{Metric} & \textbf{Formula (per agent)} \\ \midrule
Count of Orders & \F{=COUNTIF(\$B\$6:\$B\$35,A40)} \\
Total Cars Sold & \F{=SUMIF(\$B\$6:\$B\$35,A40,\$E\$6:\$E\$35)} \\
Total Sales Amount & \F{=SUMIF(\$B\$6:\$B\$35,A40,\$F\$6:\$F\$35)} \\
Average Order Amount & \F{=AVERAGEIF(\$B\$6:\$B\$35,A40,\$F\$6:\$F\$35)} \\ \midrule
Best Sales Person & \F{=INDEX(\$A\$40:\$A\$45,MATCH(MAX(\$D\$40:\$D\$45),\$D\$40:\$D\$45,0))} \\
\bottomrule\end{tabular}\end{table}

Ahmed leads with EGP 6,940,000 in total sales across 12 cars and 6 orders, ahead of
Hossam (EGP 6,770,000) despite an equal count of cars sold. The deciding factor is
the mix of car types sold, not the volume alone.

\begin{figure}[tbp]\centering
\includegraphics[width=0.92\linewidth]{fig6.png}
\caption{Task 5: per-salesperson report with the Best Sales Person highlighted and
called out below.}\end{figure}

\section{Task 6: Lookup and Calculation}

A six-model price table (Toyota, Jeep, Volkswagen, BMW, Kia, Mazda) is used to look
up the unit price for every order by car type, and the full order amount is then
calculated as price times count:

\begin{table}[htbp]\centering
\begin{tabular}{@{}ll@{}}\toprule
\textbf{Column} & \textbf{Formula} \\ \midrule
Car Price & \F{=VLOOKUP(C5,\$J\$4:\$K\$9,2,FALSE)} \\
Total Sales & \F{=D5*E5} \\
\bottomrule\end{tabular}\end{table}

\section{Task 7: Invoice}

A printable sales invoice is built on the same price-lookup principle as Task 6: the
person preparing the bill only enters the car type (from a drop-down list) and the
count; the price and the line total are calculated automatically, and the invoice
totals sum the full order:

\begin{table}[htbp]\centering
\begin{tabular}{@{}ll@{}}\toprule
\textbf{Column} & \textbf{Formula} \\ \midrule
Car Price & \F{=IF(A5="","",VLOOKUP(A5,\$H\$4:\$I\$9,2,FALSE))} \\
Total Sales & \F{=IF(A5="","",B5*C5)} \\
Total Count & \F{=SUM(C5:C15)} \\
Total Amount & \F{=SUM(D5:D15)} \\
\bottomrule\end{tabular}\end{table}

\begin{figure}[tbp]\centering
\includegraphics[width=0.82\linewidth]{fig7.png}
\caption{Task 7: a sample invoice for a three-car order, built from the same price
table.}\end{figure}

\section{Method and Validation}

\textbf{Formulas, never hard-coded results.} Every total, average, ranking and lookup
in the workbook is a live formula. No result was typed in by hand.

\textbf{Conditional formatting, not manual colouring.} The three highlighted ranges
(Tasks 2, 3 and 5) all use rule-based conditional formatting, so they stay correct if
the underlying data changes.

\textbf{One consistent style across all seven sheets.} A dark-navy header bar, EGP
currency stated once per column header, \F{dd/mm/yyyy} dates, thin borders, light row
banding, and a single font family throughout.

\textbf{Validation.} The workbook was recalculated with a full LibreOffice headless
pass across all seven sheets: 244 formulas, zero errors.

\vspace{1.4cm}
\noindent\rule{\linewidth}{0.4pt}\\[4pt]
{\small KOUAME Koffi Fid\`ele\\ Data Analysis Practice\\ koffifidelek59@gmail.com}

\end{document}
